\chapter{Grundlagen}
\label{c:grun}

\section{Techdivision \& IT-Services}
\label{s:intro:tech}
TD kurz vorstellen
Die Techdivision ist ein Unternehmen, welches 2007 in Kolbermoor, Deutschland 
gegründet wurde. Der Grundstein für die Firma wurde aber bereits 1997 in der 
Garage von Tim Wagner, dem aktuellen CTO, gelegt. Zusammen mit Thomas
Kalteis begann hier eine kleine Technologie-Schmiede. Im Jahr 2007 wurde 
daraus eine GmbH, gegründet von den Brüdern Stefan und Josef Willkommer. 


\section{Chatbots}
\label{s:intro:bot}
Laut Wikipedia ist ein Chatbot “ein Chatbot oder kurz Bot ein textbasiertes Dialogsystem, das Chatten mit einem technischen System erlaubt” \cite{Chatbot2024}. 
Es gibt verschiedene Arten von Chatbots, diese unterscheiden sich aber in der Art und Weise, wie sie auf die Eingaben des Nutzers antworten. Chatbots haben
einen Bereich zur Texteingabe und einen Bereich, in dem sie ihre Ausgabe darstellen können. Mit den meisten Chatbots kann man heutzutage über die natürliche
menschliche Sprache kommunizieren und sie versuchen auf dieselbe Art und Weise zu antworten. Mit einigen Chatbots kommuniziert man jedoch über Befehle, die
dann eine Aktion hervorrufen.\par
Die Generation des Textes eines Chatbots kann in drei verschiedene Kategorien untergeordnet werden: Retrieval based, generation based und rule based. Wir können
beobachten, wie sich diese drei Kategorien im Laufe der Zeit entwickeln.\par
Im Jahr 1950 hat Alan Turing den Turing Test vorgeschlagen, bei dem es darum geht zu erkennen, ob eine Maschine das gleiche Denkvermögen hat wie ein Mensch. Bei
diesem Test werden einem Tester zwei Gesprächspartner vorgelegt, von denen einer ein Mensch und der andere eine Maschine ist. Der Tester führt dann ein Gespräch
mit beiden Gesprächspartnern und befragt diese. Nach der Befragung wird überprüft, ob der Tester in der Lage ist, zu sagen, welcher der Partner eine Maschine ist.
Falls der Tester nicht in der Lage ist zu erkennen, welcher der Partner die Maschine ist, gilt der Turing Test als bestanden und die Maschine ist laut Turing
intelligent. Es gibt Kritik am Turing Test, welche der Meinung ist, dass der Test die Maschine nur auf ihre Funktionalität getestet und nicht wirklich auf
Intelligenz. Ein weiterer Kritikpunkt, den Menschen an dem Turing Test haben, ist, dass einige Menschen der Meinung sind, dass der Test eher die menschliche
Leichtgläubigkeit testet, als die Intelligenz der Maschine. Viele sind der Meinung, dass dieser Test die Grundsteine für die Idee der Chatbots gelegt hat.\par  
Die Geschichte der Chatbots beginnt in den 1960er Jahren mit dem ersten Chatbot aus dem Jahr 1966. Damals waren die Chatbots rule based. Der Name des Chatbots war
ELIZA und er wurde von Joseph Weizenbaum am MIT entwickelt.
Weizenbaum hat fünf Grundprobleme aufgestellt, welche er bekämpfen musste für die Entwicklung von ELIZA. Die Probleme waren:
\begin{enumerate}
    \item Die Identifikation von Stichwörtern
    \item Das Finden des minimalen Kontexts
    \item Die Auswahl der geeigneten Transformation
    \item Generation von Antworten bei Abwesenheit von Stichworten
    \item Bearbeitung der Eliza Skripte
\end{enumerate}
\cite{weizenbaumELIZAComputerProgram1966}. \par
ELIZA wurde so entwickelt, dass sie versucht hat, menschliche Sprache zu imitieren, indem sie den eingegebenen Text nach Stichwörtern durchsucht. Stichwörter sind
Wörter, denen besonders hohe Wichtigkeit laut dem ELIZA-Script gegeben wurde. Es lässt sich vermuten, dass hierbei besonders Wörter mit einem hohen Informationsgehalt
verwendet wurden. Jedes Stichwort hatte noch einen sogenannten RANK, welcher die Wichtigkeit des jeweiligen Stichworts im Vergleich zu anderen Stichworten festgelegt
hat. Nachdem die Stichwörter entdeckt wurden, geht es weiter mit einer sogenannten “Transformation Rule”. Diese besteht aus einer Zerlegungs-Phase und einer
Zusammensetzungs-Phase. Es gibt verschiedene Regeln für beides. Anhand der Syntax wird ausgemacht, welche Regel für die Zerlegung in Frage kommt und anhand dieser
wird dann die Zusammensetzung ausgemacht. /Vllt mehr dazu vllt weniger?/.
Es gibt anekdotische Hinweise von Weizenbaum und seinem Team darauf, dass es Teilnehmer gab, die eine emotionale Bindung zu seinem System aufgebaut haben. Eliza
verhält sich wie ein Psychotherapeut. Das heißt, der Nutzer beginnt das Gespräch und Eliza greift Punkte hauptsächlich aus den Eingaben des Nutzers auf, ohne großen
Zusammenhang mit Themen außerhalb dieses Gesprächs. Der Bot kennt keinen Kontext sondern antwortet, je nach seinem aktuellen Skript auf den aktuellen Satz. Selbst
Weizenbaum selbst hat scheinbar mit der Eröffnung in seine Arbeit gezeigt, dass er selbst nicht der Meinung war, dass ELIZA wirklich ein intelligentes System ist:\par
Er ist der Meinung, dass Programme im Bereich der künstlichen Intelligenz und im Bereich des heuristischen Programmierens so aufgebaut sind, dass sie auf wundersame
Weise funktionieren, sodass sie selbst erfahrene Beobachter faszinieren. Aber sobald man herausfindet, wie diese Programme funktionieren, verlieren sie ihre Magie 
und es entarnt sich als nichts mehr als eine Reihe von Anweisung, die alleine betrachtet sehr simple sind. Das Ziel seiner Forschungsarbeit ist es, die Magie hinter 
Eliza zu enttarnen, da er der Meinung ist, dass es wenig andere Dinge gibt, die diese Enttarnung so sehr nötig hatten, wie Eliza\cite{weizenbaumELIZAComputerProgram1966}.\par
Nach Eliza kam ein Chatbot namens Parry der an der Stanford University entwickelt wurde. Während Eliza eher einen Psychotherapeuten imitiert hat, hat Parry
einen Patienten mit paranoider Schizophrenie imitiert. Dieser Bot war im Jahr 1979 Teil eines Experiments, der einem Turing Test gleicht. Hier wurden verschieden
Psychotherapeuten vor Bildschirme gesetzt und mussten erraten, ob die gegenüber gestellten Aussagen von einem echten Schizophrenen Patienten kommen oder von einem
Chatbot. Teil des Experiments waren 5 Psychiater die jeweils 2 "Patienten" zur Beurteilung bekommen haben. Der erste lag bei beiden richtig, der zweite lag bei
beiden falsch. Der dritte dachte beide Patienten sind echte Patienten und die letzten zwei dachten beide sind Chatbots. Das Problem hierbei war, dass der Pool an
Probanden mit 5 Psychiatern sehr klein war. Dazu kommt noch der Fakt, dass Patienten mit Schizophrenie ein gewisses Level an Unverständlichkeit in ihren Konversationen
haben, welches die Beurteilung, ob es generell ein guter Chatbot erschwert. Heute wird Parry eher als ein schwacher Chatbot gesehen, da er nicht von Konversationen
lernen kann, keinen richtigen Kontext hat und sehr langsam ist. Natürlich ist ein Teil davon abhängig von den technischen Limitationen die es damals gab, aber es ist
dennoch wichtig zu beachten\cite{adamopoulouChatbotsHistoryTechnology2020}. 
1995 wurde der Internet Chatbot ALICE veröffentlicht. Er war im Internet zu finden und er wurde als Open Source Software veröffentlicht, hat aber nicht viele
Mitentwickler gefunden, bis die Software im Jahre 1998 von der Programmiersprache SETL auf die Programmiersprache Java gewechselt ist. Das Problem war, dass SETL
eine eher unbekannte Sprache war. Im Laufe der Jahre wurde das Projekt mehrmals mit verschiedenen Preise ausgezeichnet, wie beispielsweise dem Loebner Prize. Dieser
Bot wurde von Weizenbaums Eliza inspiriert und war selbst die Inspiration für den Film "Her", in der die Beziehung zwischen Menschen und Maschinen beleuchtet
wird\cite{moraisCanHumansFall2013}. Alice hatte eine viel größere Sammlung an Keywords, auf die sie reagieren konnte als Eliza, aber von den Grundprinzipien her,
hat sie auf die selbe Art und Weise Keywords in den Eingaben des Nutzers gesucht und dann je nach ihren Datenbankeinträgen darauf geantwortet \cite{adamopoulouChatbotsHistoryTechnology2020}.

\section{Künstliche Intelligenz}
Künstliche Intelligenz ist ein Teil der Informatik der sich damit beschäftigt, Aufgaben, die früher menschliche Intelligenz benötigt hat, von Maschinen bzw. Computern 
übernehmen zu lassen. 
Zu den Aufgaben gehören unter anderem:
\begin{itemize}
    \item Erkennen von Mustern
    \item Treffen von Entscheidungen
    \item Verstehen von Sprache
\end{itemize}
Durch Fortschritte in der künstlichen Intelligenz konnten Forstschritte in all diesen Aufgabenbereichen erreicht werden.
Die verschiedenen Aufgaben werden teils in verschiedenen Teilbereichen der künstlichen Intelligenz behandelt. Ein paar der Teilbereiche sind:
\begin{itemize}
    \item Computer Vision
    \item Natural Language Processing
    \item Robotik
\end{itemize}


\subsection{Machinelles Lernen}


\subsection{Neuronale Netzwerke}
\label{s:intro:nn}
Neuronale Netze sind ein Teil des Maschinellen Lernens, welche versucht sich stark an dem Aufbau eines biologischen Gehirns zu orientieren.
Dementsprechend bestehen sie aus Tausenden oder sogar Millionen ``Neuronen'', welche durch Knotenpunkte dargestellt werden. Diese ``Neuronen''
sind stark miteinander verbunden. Sie sind über verschiedene Reihen miteinander verbunden, wie in der Abbildung 2.1 dargestellt.\par

Jeder Eingang an einem Neuron hat eine Gewichtung. Wenn die Gewichtung eines Eingangs 0 ist, dann heißt es, dass zwei Knoten 
keine Verbindung miteinander haben. Wenn die Gewichtung positiv ist, dann hat es eine bestärkende Wirkung und, falls die Gewichtung negativ
ist, dann hat es eine hemmende Wirkung. \par %wikipedia https://de.wikipedia.org/wiki/Künstliches_Neuron
Jedes Neuron hat zudem einen sogenannten Bias, der Bias wird benötigt, damit man die Aktivierungsschwelle der Aktivierungsfunktion 
verschieben kann. Dies sorgt für eine erhöhte Flexibilität, da der Wert der Aktivierungsfunktion ohne den Bias immer null wäre, wenn alle
Eingaben null sind. Durch den Bias hat das Neuron also eine Grundaktivität, welche während dem Training angepasst werden kann, damit die Ergebnisse 
besser für den eigenen Fall werden. 
Das aktuelle Neuron summiert dann alle gewichteten Werte und addiert den Bias. Wie in dieser Formel dargestellt:
\[
\sum_{i} w_i x_i + b
\] 
Der summierte Wert wird durch die Aktivierungsfunktion transformiert und das Ergebnis wird dann dem nächsten Neuron in der nächsten Schicht weitergeleitet.
Typische Aktivierungsfunktionen sind
\begin{itemize}
    \item Sigmoid Function : \[f(x) = \frac{1}{1 + e^{-x}}\]
    \item ReLu (Rectified Linear Unit): \[f(x) = \max(0, x)\]
    \item Tanh: \[f(x)=tanh(x)\]
\end{itemize}
\cite{kempkesNeuronUndAktivierungsfunktion2023}

\begin{figure}[t]
    \centering
    \includegraphics[width=\linewidth]{\figdir/neuron}
    \caption{Neuron eines Neuronalen Netzwerks}
\end{figure}
Die zwei Schichten die es immer gibt sind die Eingabe Schicht und die Ausgabe Schicht. Dazwischen kann es verstecke Schicht geben, die benötigt werden um 
komplexere Zusammenhänge zu bearbeiten, muss es aber nicht.
Die Eingabe Schicht erhält die Eingabedaten von außen und gibt diese dann weiter an die nächste Schicht. 
Als nächstes werden die Daten weiter an eine verborgene Schicht gegeben, die die Werte dann bearbeitet. 
Es gibt X verschieden verborgene Reih in einem Neuronalen Netzwerk. Die Anzahl der Reihen und die Anzahl 
der Neuronen pro versteckter Schicht kann sich unterscheiden. Es gibt hierbei keine allgemeingültige Größe,
die für alle Fälle passt. Stattdessen kommen hierbei Erfahrung und Experimentation ins Spiel.\par 
Es kann sogar sein, dass gar keine versteckte Reihen notwendig sind, wenn das Model linear seperierbar ist. 
Das liegt daran, dass ein Modell ohne versteckte Schichten bereits in der Lage ist diese Grenze zu erlernen 
und die Daten dann korrekt zu klassifizieren.\par
Die Anzahl der Neuronen schwankt zwischen verschiedenen Neuronalen Netzen und auch innerhalb eines Neuronalen Netzes schwankt 
die Anzahl der Neuronen zwischen den verschiedenen Schichten.
Bei der Eingabe Schicht kommt es auf die Trainingsdaten an, diese legen die Anzahl der Neuronen fest. 
Um genau zu sein, kommt es auf die Anzahl der Features an, die die Trainingsdaten beinhalten. Bei einem 
Neuronalen Netz, welches mit Daten aus einer normalen Tabelle trainiert wird, ist das beispielsweise die Anzahl
der Spalten, bei einem bei einem Neuronalen Netz, welches mit Bildern trainiert werden soll, könnten das die 
einzelnen Pixel des jeweiligen Bildes sein.\par
\cite{krishnanHowDetermineNumber2022}
Die Anzahl der Neuronen in der Ausgabe Schicht kommt darauf an, welche Art von Modell man betrachtet.
Es wird hier vor allem zwischen zwei verschiedenen Arten unterschieden. Die zwei Arten die es gibt sind 
Regressoren und Klassifikatoren. Bei Regressoren gibt es ein Ausgabe Neuron. Bei Klassifikatoren kann es 
mehrere geben.\par Regressoren ordnen die Ergebnisse im Regelfall in eine von zwei Optionen ein. Der Wert den
das Modell berechnet ist ein Wert zwischen 0 und 1. Aber die Ausgabe des Neuronalen Netzes ist entweder
eins oder null. Hierbei gibt es einen Schwellenwert, der, wenn er überschritten wird, die Ausgabe des Modells 
auf eins bringt und, wenn dieser Wert nicht überschritten wird, ist die Ausgabe null. Ein Regressor 
klassifiziert die Daten in eine von zwei Gruppen, wenn der Wert null ist, gehört die Eingabe in eine Gruppe
und bei eins in die jeweils andere.\par \par
Beispiele hierfür sind: 
\begin{itemize}
    \item Gesichtserkennung
    \item Stimmerkennung
    \item Spamfilter in emails
\end{itemize}
Die Klassifikatoren haben ein oder auch mehrere Ausgabe Neuronen. Die Anzahl der Neuronen kommt auf die Anzahl der möglichen Ausgaben an.
Beispielsweise könnte man ein System betrachten, welches lernen soll, welche Art von Fahrzeug auf einem Bild dargestellt ist. 
In der Abbildung 2.2 ist ein simpler Decision Tree dargestellt. Anhand dieser kann man sehen, welche Optionen es als Asugaben gibt. 
In dieser Situation hätte die Ausgangs Schicht drei Neuronen, da es drei mögliche Ausgaben gibt. \par

%vmtl raus nehmen is kinda dogshit
\begin{figure}[t]
    \centering
    \includegraphics[width=\linewidth]{\figdir/decision-tree}
    \caption{Neuronale Netz}
\end{figure}

%hier definitiv neues bild machen das kleiner und besser ist 
\begin{figure}[t]
    \centering
    \includegraphics[width=\linewidth]{\figdir/nn}
    \caption{Neuronale Netz}
\end{figure}
\subsection{Training}
Um ein Modell nach der 

\section{LLMs}
\label{s:intro:llm}
Large Language Models sind eine Art der künstlichen Intelligenz, welche mithilfe probalistischer Methoden Text generiert. 
Im Grunde sind diese LLMs Text Vervollständigungsgeneratoren. Heutzutage basieren Chatbots auf sogenannten Large Language 
Models. LLMs sind in ihrem Kern nur 2 Dateien. Eine Datei die die sogenannten Gewichte des Neural Networks beinhaltet, 
welche beim Trainieren entstanden sind und eine Datei, welche den Code beinhaltet,damit man dieses Neural Network ausführen kann.
Der Code lässt sich in ein paar hundert Zeilen Code schreiben, selbst ohne Abhängigkeiten. Der größere und aufwändigere Part sind 
die Gewichte. Die Gewichte entstehen, beim trainieren des Neural Networks, welches auf der Transformer - Decoder Architektur 
basiert, welche 2017 im Paper "All you need is attention" vorgestellt wurde. Der genaue Stack den bspw. OpenAI heutzutage benutzt 
ist nicht bekannt, aber wir wissen, dass dieser ebenfalls auf der Transformer Architektur basiert. 
\begin{figure}[htbp]
    \centering
    \begin{tikzpicture}
        \begin{axis}[
            xlabel=Woche,
            ylabel=LLM,
            date coordinates in=x,
            xticklabel style={rotate=45,anchor=near xticklabel},
            date ZERO=2019-05-12, % Set the reference date
            ]
        
            \addplot table [x=Woche, y=LLM, col sep=comma] {data/LLM.csv};
        
        \end{axis}
    \end{tikzpicture}
    \caption{LLM Google Trends}
    \label{fig:trend_llm}
\end{figure}
Wir sehen anhand der \ref{fig:trend_llm}, welche über Google Trends das Interesse am Wort "LLM" im Verlauf der letzten 5 Jahre darstellt, dass das Interesse
an Large Language Models besonder stark im Dezember 2022 angestiegen ist und seitdem immer weiter steigt. 
Die Daten von Google Trends sind normalisiert zwischen 0 und 100 also können wird dadurch zwar keine genauen Angaben machen, wir sehen aber, dass das
Wachstum stetig am steigen ist.\par
Im Grunde sind die LLMs die hinter Chatbots wie ChatGPT von OpenAI und LeChat von MistralAI sogenannte GPTs. Wenn man den Begriff aufteilt sieht man
was dahinter steckt. GPT steht für Generative pretrained Transformer. Generative steht für die Art und Weise
auf der diese LLM ihre Ausgaben generiert, Pretrained steht dafür, dass die LLM vor ihrer Nutzung auf einer Datenmenge trainiert wurde und Transformer
steht hier für die Transformer Architektur.

\subsection{Tokenizer}
\label{s:intro:llm:token}
Tokenizer sind ein wichtiger Part der Machine Learning Toolchain 
im Zusammenhang mit dem Lernen von menschlicher Sprache.
Wenn man im Bereich des Natural Language Processing ein Machine 
Learning trainieren möchte stellt sich die Frage, wie man Texte
aufteilen kann, damit das Modell, welches trainiert wird diese 
verarbeiten kann.\par
Die Tokenizer teilen die Texte so auf, dass das zu trainierende
Modell komplexe Sätze besser aufteilen und verstehen kann. 
Der Tokenizer macht das, indem er den Text in sogenannte
Tokens aufteilt. Ein Token kann ein ganzes Wort, ein Teil eines
Worts oder ein Satzzeichen sein.\par
Dies ist ein wichtiger Teil des \ac{nlp}, da die Aufteilung der 
Trainingsdaten einen großen Einfluss darauf hat, wie gut die 
\ac{llm} funktioniert. Genau wie bei Compilern, die Programme in
Tokens aufteilen, benutzen wir bei \ac{llm} Tokenizer, damit wir
die Texte aufteilen können und die dabei entstandenen Tokens über 
verschiedene grammatikalische Regeln zusammensetzen können.\par
Eines der größten Probleme hierbei ist, dass verschiedene Sprachen verschiedene
grammatikalische Regeln haben und somit verschiedene Arten des 
Zerlegens oder Zusammensetzens brauchen. 
Ein weiteres Problem ist, dass das selbe Wort teils verschiedene Bedeutungen haben
kann. Hierbei ist der Kontext wichtig und die Anzahl der Tokens, die das 
Modell erkennen kann. Man könnte hierbei alle möglichen Token aufnehmen, 
was aber dafür sorgen kann, dass das System overfitted beim Training.
Dazu kommt, dass das die notwendigen Kosten zum Trainieren und für die Inference 
des Systems erhöhen würde.\cite{perryWhatTokenizationNatural2021} \cite{chenStrategiesTrainingLarge2016}


\subsection{Training}
\label{s:intro:llm:train}
Der erste Schritt um eine \ac{llm} zu erstellen ist es diese zu trainieren. 
Die \ac{llm} trainieren heißt in diesem Fall die LLM mit Daten zu trainieren.
Dafür werden viele verschiedene Daten gesammelt und die LLM wird dann 
mit Hilfe dieser Daten trainiert.\par
Die \ac{llm} teilt in diesem Prozess die Daten die sie erhalten hat mit Hilfe eines
Tokenizers in verschieden Tokens ab und lernt in diesem Zusammenhang, den 
Aufbau menschlicher Sprache in der jeweiligen Sprache, die zum trainieren 
verwendet wird. \par

\subsection{Hypertraining}


\subsection{Inference}
\label{s:intro:llm:inf}
Sobald ein Model trainiert wurde und bereit ist benutzt zu werden, muss man es ausführen. Dieser Prozess nennt sich Inference.
Um Inference ausführen zu können gibt es heutzutage verschiedene Möglichkeiten. Es gibt zwar die Möglichkeit Inference über 
online APIs wie beispielsweise die Huggingface Inference API \cite{ServerlessInferenceAPI} oder ähnliche Services auszuführen, 
wir behandeln hier aber nur lokale Optionen. \par
Das Hauptproblem bei der Inference von diesen Modellen war zu Beginn, dass sie sehr viele Ressourcen benötigt haben und nur über
spezielle Hardware ausführbar waren. Spezielle Hardware heißt in diesem Fall, dass sie zu Beginn nur über Grafikkarten mit sehr
viel VRAM nutzbar waren. 

\subsection{Prompting}
\label{s:intro:llm:prom}

\subsection{Encoder-Decoder Architektur}
Bevor auf die Transformer Architektur eingegangen werden kann, muss auf die vorangeganen Architekturen eingeganen werden.
%erklaeren wieso das erklaert werden muss
Diese Technologien haben sich damals vorallem im Kontext des Sequence-to-Sequence Learnings weiterentwickelt. 
Sequence-to-Sequence Learning oder kurz Seq2Seq Learning, bringt eine Sequenz variabler Länge in eine andere
Sequenz variabler Länge.\cite{sutskeverSequenceSequenceLearning2014} Beispiele hierfür wären Text Zusammenfassungen und Übersetzungen. Die Architektur aus der sich
der Transformer, der in dem Abschnitt \ref{s:intro:llm:trans} behandelt wird, der in den LLMs heutzutage 
verwendet wird, entwickelt hat, ist die Encoder - Decoder Architektur.\par
Die Encoder-Decoder Architektur besteht aus zwei Rekurrenten Neuralen Netzen. Wie ein Neurales Netz funktioniert wurde bereits im 
Abschnitt \ref*{s:intro:nn} behandelt. Das besondere an einem Rekurrenten Neuralen Netz ist, dass dieses die vorherigen Aktivierungen 
der einzelnen Neuronen mit berücksichtigt. Dadurch, dass die vorangeganen Aktivierungen ebenfalls beachtet werden, können Kontexte 
erkannt werden und Sätze können sinnvoll übersetzt werden.\par
Hierbei erhählt der Encoder eine Eingabe Sequenz, welche dann Kontext Vektoren erstellt, welcher vom Decoder in eine Ausgabe Sequez
umgewandelt wird.  
Das Problem bei der originalen Encoder-Decoder Architektur ist, dass die Vektoren eine festgesetzte Größe hatten und  


\section{Transformer}
\label{s:intro:llm:trans}
Die Technologie auf der die heutigen LLMs im Grunde basieren ist die sogenannte Transformer Technologie.\par
Diese Architektur wurde im Research Paper "Attention is all you need" im Jahr 2017 von Google Forschern vorgestellt. Laut ihrer Forschung behaupteten 
sie, dass ihre neue Architektur signifikant weniger Zeit zum Training braucht, während es besser parallelisierbar ist und die Ergebnisse trotzdem
besser sind. Nach 3.5 Tagen Training auf nur 8 Grafikkarten hat ihr Modell einen neuen Single Model Rekord aufgestellt 
in der WMT 2014 Englisch-zu-Französisch Übersetzungsaufgabe\cite{vaswaniAttentionAllYou2023}. 
\begin{figure}[t]
\centering
\includegraphics[width=\linewidth]{\figdir/transformer}
\caption{Transformer Architektur}
\end{figure}
Von Yuening Jia - DOI:10.1088/1742-6596/1314/1/012186, CC BY-SA 3.0, https://commons.wikimedia.org/w/index.php?curid=121340680

\subsection{Attention Mechanismus}
Die Grundidee auf die die Transformer aufbauen ist die des Attention Mechanism, welcher 2014 im Paper ``Neural Machine Translation by Jointly Learning
to Align and Translate'' von Bahdanau et al. vorgestellt wurde.
Beim Attention Mechanism geht es darum, Deep Learning Modelle zu verbessern, indem 
Der Attention Mechanism basiert 
\cite{bahdanauNeuralMachineTranslation2016}

\section{Vektor Datenbanken}
\label{s:intro:dat}
Vektor Datenkbanken sind die Datenbanken.

\section{Embeddings}
\label{s:intr:emb}


\section{RAG}
\label{s:intro:rag}
Eines der größten Probleme der LLMs heutzutage ist, dass die Informationen, die sie zum beantworten für Fragen benutzen, 
teilweise veraltet sind oder sehr allgemein. Ein weiteres großes Problem ist die im Abschnitt \ref{s:intro:llm} genannte 
Halluzination, die passieren kann, wenn eine LLM eine Frage zu der es eigentlich nichts weiß, beantwortet. Eine Technologie, 
die versucht, dieses Problem zu beseitigen, ist die sogenannte Retrieval Augmented Generation. 
Die Retrieval Augmented Generation soll das parametrische Gedächtnis der LLM
mittels eines nicht parametrischen Gedächtnisses erweitern.\cite{lewisRetrievalAugmentedGenerationKnowledgeIntensive2021} \par 
Retrieval Augmented Generation löst ein paar Probleme vor denen LLMs heutzutage 
stehen, ist aber keine perfekte Lösung und muss sich auch seinen eigenen Problemen
stellen. In diesem Zusammenhang gibt kann man heute RAG in drei Stufen 
einordnen.\cite{gaoRetrievalAugmentedGenerationLarge2024}

\subsection*{Naive RAG}
%schreib das nochmal um das ist kinda hart geklaut
Naive RAG ist die traditionelle Herangehensweise and das Problem der RAG.
Sie besteht aus drei Stufen. Indexing(Indizierung??), Abfrage und Generierung.\par
Beim Indizieren werden die rohen Eingabedaten gereinigt und extrahiert.
Dies ist wichtig, da es verschieden Formate haben kann, bei Textdaten 
gibt es bspw. PDF Dateien, Markdown oder HTML. Diese Texte werden nach der Reinigung
in Plain Text umgewandelt und dieser Plain Text wird dann in kleinere Teile unterteilt .
Diese kleineren Parts werden dann über ein Embedding modell in eine 
Vektordarstellung umgewandelt und dann in einer Vektordatenbank abgespeichert.\par
Bei der Abfrage erhält das RAG System eine Anfrage und benutzt das selbe 
Embedding Modell, wie in der Indizierungsphase, um die Abfrage in eine 
Vektordarstellung umzuwandeln. Im Anschluss wird durch einen Semantischen
Vergleich nach den Top-k Ergebnissen gesucht und diese werden dann 
ausgegeben.\par
Bei der Generation bekommt die LLM dann die Abfrage und die ausgegeben 
Ergebnisse in einem Prompt zurück und generiert daraus dann eine Antwort. 
Was genau hier beachtet wird kommt auf die LLM, das Kontext Fenster und den
 benutzten Prompt an\cite{gaoRetrievalAugmentedGenerationLarge2024}.\par
Wie in \ref{s:intro:rag} bereits erwähnt gibt es hierbei ein paar Probleme.
Das Hauptproblem, welches hier besteht ist die der Abfrage. Bei der Abfrage 
gibt es oft Probleme, welche mit der Genauigkeit zu tun haben. Dementsprechend
kann es passieren, dass das RAG System falsche oder unpassende Informationen
zurück gibt. Aufgrund der Probleme bei der Abfrage entstehen dann ebenfalls 
Probleme bei der Generation der Antwort. Was dazu führt, dass die Qualität 
der Antworten sinkt. Dadurch sinkt dann ebenfalls die Verlässlichkeit auf 
das System. 

\subsection{Advanced RAG}
Advanced RAG wurde entwickelt um den Problemen der naiven RAG entgegenzuwirken
bzw. sie zu beseitigen. Dementsprechend werden hier einige Dinge geändert. Vorallem
in der Stufe der Abfrage wurde hier einiges geändert. Um die Qualität zu erhöhen 
wurden der Abfrage pre- und post Strategien hinzugefügt.
Dazu wurde der Indexing Prozess angepasst.
Beim Indexing wird nun Metadata mit eingebunden und der Text wird genauer unterteilt.
Pre-retrieval:
Beim pre-retrieval wird vorallem darauf geachtet die Anfrage so klar wie möglich zu machen.
Hierfür werden Metadaten angehängt, Anfragen umgewandelt, Anfragen erweitert und noch 
einige weitere Dinge abgeändert. Das Ziel hier ist es die Anfrage zu ändern, sodass sie so 
klar wie möglich ist und sie für Retrival zu optimieren.\par
Post Retrieval: 
Nachdem der Kontext nach dem Retrieval Prozess abgefragt wurde, wird
versucht die abgerufenen Daten möglichst optimal mit der Anfrage zu verbinden.
Hierfür werden vorallem das komprimieren und das neu Anordnen des Kontexts 
verwendet. Die wichtigsten Dokumente werden an eines der Enden der Anfrage 
geschoben und es wird versucht die wichtigsten Informationen heraus zu suchen
und nur diese zurück zu geben. Dadurch sollte die LLM genug Informationen 
erlangen, aber nicht zu viele und die Antworten sollten besser werden.

\subsection{Modular RAG}

\section{Ragas}
\label{s:intro:ragas}

\section{Langchain}
\label{s:intro:lang}

\section{Huggingface}
\label{s:intro:hug}

\section{llama.cpp \& Ollama}
\label{s:intro:oll}
Ein Problem, welches 

\section{Google}
\subsection{Google Cloud}
Dadurch, dass die Techdivision ein Google Partner ist und intern versucht 
möglichst viele Produkte von Google zu nutzen 

\subsection{Google Chat}
Google Chat ist eine von Google entwickelte Chatplattform. Diese Plattform 
hat 2017 Google Hangout zusammen mit Google Meet ersetzt. Zu Beginn war das 
Projekt nur im Zusammenhang mit Google Workspace erhältlich und somit war es 
eine Chatplattform die hauptsächlich für Teams und Unternehmen benutzbar war. 
2021 wurde die Plattform im Early Access auch für normale Benutzer erreichbar 
gemacht und ist seit 2022 für jeden benutzbar.\par
Diese Plattform ist die Plattform, die die Techdivision hauptsächlich intern 
zur Kommunikation benutzt. Es gibt auf der einen Seite 1 zu 1 Chat in der eine 
Person mit nur einer anderer Person kommuniziert und auf der anderen Seite gibt es 
Innerhalb dieses Services die Möglichkeit 
sogenannte Spaces zu erstellen, welche Gruppenchats ähneln. Ein Teil dieser
Gruppenchats sind die in \ref*{s:intro:tech} genannten Supportchats in die 
der Chatbot final eingebaut werden soll.